% SI Text 1: derivation of the gap rule for constant-returns production functions (Proposition 1, Corollaries 1 and 2, Examples 1 and 2, Fig. S1).
\section{Derivation of the gap rule}\label{app:gap-rule}

\subsection*{The within-round allocation problem}

A population of $n$ researchers. Researcher $i$ has capability $K_i > 0$ and baseline resources $R_{i0} \geq 0$, both fixed within the round. A grant $g_i \geq 0$ adds to researcher $i$'s resources for the round, so that their expected output is $F(K_i,\, R_{i0} + g_i)$. We prove the gap rule for every production function of the form
\[
F(K, R) = A\,K\,h(R/K),
\]
where $A > 0$ is the productivity constant and $h\colon [0, \infty) \to [0, \infty)$ is continuous, continuously differentiable on $(0, \infty)$, strictly increasing, and strictly concave, so that $h'$ is positive and strictly decreasing. This is the class of production functions with constant returns to scale in $(K, R)$ and diminishing returns to resources. The harmonic mean of the main text is the instance $h(x) = x/(1 + x)$, for which $F(K, R) = AKR/(K + R)$; every member of the power-mean family of \ref{app:technology} with exponent $\gamma < 1$ is an instance, with $h(x) = \big((1 + x^{\gamma})/2\big)^{1/\gamma}$. Write
\[
\varphi_i(g) = A\,K_i\,h\!\left(\frac{R_{i0} + g}{K_i}\right)
\]
for researcher $i$'s expected output as a function of their grant alone; for the harmonic mean, $\varphi_i(g) = A K_i (R_{i0} + g)/(K_i + R_{i0} + g)$.

Throughout, the round's budget is $B > 0$. Nothing else about the model is used, so the result holds for every realized population and every round's budget; across rounds, the constant $c$ and the funded set are recomputed from the new population state.

The funder observes $(K_i, R_{i0})_{i=1}^n$ and divides the budget so as to maximize expected output:
\[
(\mathrm{P}) \qquad \max_{g_1, \ldots, g_n \geq 0} \; \sum_{i=1}^n \varphi_i(g_i) \quad \text{subject to} \quad \sum_{i=1}^n g_i \leq B.
\]

\subsection*{Three lemmas}

\begin{lemma}[the return to a grant]\label{lem:return}
For each $i$, on $g \in [0, \infty)$: (i) $\varphi_i$ is strictly increasing, with $\varphi_i'(g) = A\,h'\!\big((R_{i0} + g)/K_i\big) > 0$; (ii) $\varphi_i$ is strictly concave, since $\varphi_i'$ is strictly decreasing in $g$; (iii) $\varphi_i(g) \uparrow A K_i \bar h$ as $g \to \infty$, where $\bar h = \lim_{x \to \infty} h(x) \leq \infty$. For the harmonic mean, $\varphi_i'(g) = A\,K_i^2/(K_i + R_{i0} + g)^2$, $\varphi_i''(g) = -\,2A\,K_i^2/(K_i + R_{i0} + g)^3$, $\varphi_i'(g) < A$ whenever $R_{i0} + g > 0$, and $\bar h = 1$, so expected output never exceeds $AK_i$.
\end{lemma}

\begin{proof}
Differentiating $\varphi_i(g) = A K_i h((R_{i0} + g)/K_i)$ in $g$ gives (i), and (ii) follows because $h'$ is strictly decreasing and its argument increases with $g$. For (iii), $h$ is increasing, so $h(x) \uparrow \bar h$ as $x \to \infty$. For the harmonic mean, write $S = K_i + R_{i0} + g$; then $\varphi_i(g) = AK_i - AK_i^2/S$, and the stated derivatives follow (note $dS/dg = 1$); $\varphi_i'(g) = AK_i^2/S^2 < A$ exactly when $S > K_i$, that is, when $R_{i0} + g > 0$.
\end{proof}

\Cref{lem:return} says that the marginal value of a dollar to researcher $i$ depends on capability and resources only through their ratio, decreases as that ratio rises, and so increases with capability and decreases with the resources they already hold; that each dollar granted lowers the value of the next; and that, when $h$ is bounded, expected output never exceeds a ceiling set by capability.

\begin{lemma}[existence, uniqueness, exhaustion]\label{lem:existence}
Problem $(\mathrm{P})$ has exactly one solution $g^* = (g_1^*, \ldots, g_n^*)$, and it exhausts the budget: $\sum_i g_i^* = B$.
\end{lemma}
\begin{proof}
The feasible set is compact and convex and the objective $f(g) = \sum_i \varphi_i(g_i)$ is continuous and, by \cref{lem:return}(ii), strictly concave, so a maximizer exists and is unique. If $\sum_i g_i^* < B$, increasing any one $g_i^*$ slightly stays feasible and raises $f$ (\cref{lem:return}(i)), contradicting optimality.
\end{proof}
\begin{lemma}[equal marginal value]\label{lem:emv}
A feasible allocation $g$ with $\sum_i g_i = B$ solves $(\mathrm{P})$ if and only if there is a $\nu > 0$ such that
\[
\varphi_i'(g_i) = \nu \;\text{ for every } i \text{ with } g_i > 0, \qquad \varphi_i'(0) \leq \nu \;\text{ for every } i \text{ with } g_i = 0,
\]
where $\varphi_i'(0)$ may be $+\infty$ (as it is when $R_{i0} = 0$ and $h'(0^+) = \infty$), in which case researcher $i$ is funded.
\end{lemma}

\begin{proof}
\emph{Necessity.} Let $g^*$ solve $(\mathrm{P})$. Since $B > 0$ is exhausted (\cref{lem:existence}), some researcher is funded: $g_i^* > 0$ for some $i$. Fix any such $i$ and any $j \neq i$, and for $0 < \delta \leq g_i^*$ consider the transfer that replaces $g_i^*$ by $g_i^* - \delta$ and $g_j^*$ by $g_j^* + \delta$, leaving the rest unchanged. The result is feasible, so optimality gives
\[
\varphi_j(g_j^* + \delta) - \varphi_j(g_j^*) \;\leq\; \varphi_i(g_i^*) - \varphi_i(g_i^* - \delta).
\]
Dividing by $\delta$ and letting $\delta \downarrow 0$ yields $\varphi_j'(g_j^*) \leq \varphi_i'(g_i^*)$. If $j$ is also funded, the same argument with $i$ and $j$ exchanged gives the reverse inequality, so all funded researchers share a common marginal value; call it $\nu = \varphi_i'(g_i^*) > 0$ (positivity by \cref{lem:return}(i)). For unfunded $j$, the displayed inequality reads $\varphi_j'(0) \leq \nu$.

\emph{Sufficiency.} Let $g$ satisfy the condition with multiplier $\nu$, and let $\tilde g$ be any feasible allocation. Concavity of $\varphi_i$ gives, for every $i$, $\varphi_i(\tilde g_i) \leq \varphi_i(g_i) + \varphi_i'(g_i)\,(\tilde g_i - g_i)$. For funded $i$, $\varphi_i'(g_i) = \nu$, so $\varphi_i'(g_i)(\tilde g_i - g_i) = \nu\,(\tilde g_i - g_i)$. For unfunded $i$, $g_i = 0$ and $\tilde g_i - g_i = \tilde g_i \geq 0$, so $\varphi_i'(0)\,(\tilde g_i - g_i) \leq \nu\,(\tilde g_i - g_i)$. Summing over $i$:
\[
f(\tilde g) \;\leq\; f(g) + \nu \sum_i (\tilde g_i - g_i) \;=\; f(g) + \nu\Big(\sum_i \tilde g_i - B\Big) \;\leq\; f(g),
\]
since $\tilde g$ is feasible and $\nu > 0$. So $g$ solves $(\mathrm{P})$.
\end{proof}

\Cref{lem:emv} is the formal content of the sentence in the main text: an allocation is optimal exactly when no dollar can be moved to a researcher who can produce more value with it.

\subsection*{The proposition}

\begin{proprestate}[the gap rule]
Let $F(K, R) = A\,K\,h(R/K)$ with $h$ increasing and strictly concave. Then problem $(\mathrm{P})$ has the unique solution
\[
g_i^* = \max(c\,K_i - R_{i0},\, 0), \qquad i = 1, \ldots, n,
\]
where $c > 0$ is the unique constant satisfying the budget identity $\sum_{i=1}^n \max(c\,K_i - R_{i0},\, 0) = B$. Equivalently, $c = (h')^{-1}(\nu/A)$, where $\nu$ is the common marginal value of \cref{lem:emv}; for the harmonic mean, $c = \sqrt{A/\nu} - 1$ with $\nu \in (0, A)$.
\end{proprestate}

\begin{proof}
Let $g^*$ be the unique solution (\cref{lem:existence}) and $\nu$ the multiplier that \cref{lem:emv} associates with it. The proof proceeds in four steps.

\emph{Step 1: $c = (h')^{-1}(\nu/A)$ is well defined and positive.} Some researcher $i$ is funded, and for them \cref{lem:emv} gives $\nu = \varphi_i'(g_i^*) = A\,h'\big((R_{i0} + g_i^*)/K_i\big)$, so $\nu/A$ lies in the range of $h'$; since $h'$ is strictly decreasing, its inverse is defined there, and $c = (h')^{-1}(\nu/A) = (R_{i0} + g_i^*)/K_i > 0$ because $g_i^* > 0$. For the harmonic mean, $h'(x) = 1/(1 + x)^2$, so $(h')^{-1}(y) = 1/\sqrt{y} - 1$ and $c = \sqrt{A/\nu} - 1$; positivity of $c$ is then $\nu < A$, which \cref{lem:return} gives directly.

\emph{Step 2: funded researchers receive exactly their gap.} Let $g_i^* > 0$. By \cref{lem:emv}, $A\,h'\big((R_{i0} + g_i^*)/K_i\big) = \nu$, so $(R_{i0} + g_i^*)/K_i = (h')^{-1}(\nu/A) = c$ and $g_i^* = c\,K_i - R_{i0}$. Since $g_i^* > 0$, this equals $\max(c\,K_i - R_{i0},\, 0)$. Every funded researcher is thus brought to the same ratio of resources to capability, $c$.

\emph{Step 3: unfunded researchers have no gap.} Let $g_i^* = 0$. By \cref{lem:emv}, $A\,h'(R_{i0}/K_i) = \varphi_i'(0) \leq \nu = A\,h'(c)$, and since $h'$ is strictly decreasing, $R_{i0}/K_i \geq c$, that is, $R_{i0} \geq c\,K_i$. Then $\max(c\,K_i - R_{i0},\, 0) = 0 = g_i^*$.

Steps 2 and 3 establish $g_i^* = \max(c\,K_i - R_{i0}, 0)$ for every $i$, and the budget identity is exhaustion (\cref{lem:existence}). It remains to show $c$ is the only constant satisfying that identity.

\emph{Step 4: the budget identity pins $c$ uniquely.} Define $G(c) = \sum_{i=1}^n \max(c\,K_i - R_{i0},\, 0)$ for $c \geq 0$. $G$ is continuous and nondecreasing, with $G(0) = 0$ and $G(c) \to \infty$ as $c \to \infty$. Moreover $G$ is strictly increasing wherever it is positive: if $G(c) > 0$, then $c\,K_j > R_{j0}$ for some $j$, and for any $c' > c$, $G(c') \geq G(c) + K_j\,(c' - c) > G(c)$. Now suppose $c_1 < c_2$ both satisfy $G(c) = B$. Since $B > 0$, $G(c_1) > 0$, so $G(c_2) > G(c_1) = B$, a contradiction.
\end{proof}

The proof uses only that the marginal value of a dollar depends on $K$ and $R$ through the ratio $R/K$ and falls as that ratio rises. Constant returns to scale give the first property and diminishing returns to resources the second; the harmonic form contributes nothing beyond the closed form of $c$. The constant $c$ is the ratio of resources to capability at which the optimal allocation leaves every funded researcher, and so also the rate at which the rule converts a unit of capability into a target level of resources.

\subsection*{Consequences}

\begin{corollary}[funding frontier and budget comparative statics]\label{cor:frontier}
(i) Researcher $i$ is funded, $g_i^* > 0$, exactly when $K_i > R_{i0}/c$. (ii) As a function of the budget $B$, the constant $c(B)$ is continuous and strictly increasing; hence the funded set is nondecreasing in $B$, every funded researcher's grant $c(B)\,K_i - R_{i0}$ is strictly increasing in $B$, and the common marginal value $\nu(B) = A\,h'(c(B))$ is strictly decreasing in $B$. Both parts follow from the formula and Step 4.
\end{corollary}
\begin{corollary}[the value of targeting vanishes in the budget]\label{cor:targeting-vanishes}
Let $f^*(B)$ denote the optimal value of $(\mathrm{P})$ and $f^u(B) = \sum_i \varphi_i(B/n)$ the value of uniform funding. Then $0 \leq f^*(B) - f^u(B)$ for every $B$. If $h$ is bounded, $\bar h < \infty$, then $f^*(B) - f^u(B) \to 0$ as $B \to \infty$.
\end{corollary}

\begin{proof}
The first inequality holds because uniform funding is feasible. For the second: by \cref{lem:return}(iii), $\varphi_i(g) < A K_i \bar h$ for every $g$, so $f^*(B) \leq A\bar h\sum_i K_i$; and $f^u(B) = \sum_i \varphi_i(B/n) \to A\bar h\sum_i K_i$ as $B \to \infty$. Hence $0 \leq f^*(B) - f^u(B) \leq A\bar h\sum_i K_i - f^u(B) \to 0$.
\end{proof}

Two remarks. First, the corollary says nothing about monotonicity: $f^*(B) - f^u(B)$ is zero at $B = 0$ and positive for small $B > 0$ whenever researchers differ, so it cannot decrease throughout. The decrease reported in the main text is a numerical finding about the ratio $(f^* - f^u)/(f^u - f^0)$, with $f^0$ the output under no funding: it fell monotonically in the budget, on budget scales from $0.02$ to $3$, in every one of 300 populations of 50 researchers drawn from the default distributions, and in the population of Fig.~\ref{fig:targeting-value}. Since uniform funding's own gain over no funding does not vanish (it approaches $A\bar h\sum_i K_i - \sum_i \varphi_i(0) > 0$), the ratio vanishes with the difference. Second, boundedness of $h$ is what makes targeting's value vanish: in the power-mean family, $h$ is bounded exactly when $\gamma < 0$, with $\bar h = 2^{-1/\gamma}$ (so $\bar h = 1$ for the harmonic mean), whereas under Cobb-Douglas production ($\gamma \to 0$, $h(x) = \sqrt{x}$) output has no ceiling set by capability and $f^*(B) - f^u(B)$ grows without bound in $B$ for any population with unequal capabilities. The Leontief form $h(x) = \min(1, x)$ is the limit $\gamma \to -\infty$ of the family; it is not strictly concave, so the proposition does not apply to it directly, and it is treated numerically in \ref{app:technology}.

\begin{example}[funding the track record can produce less than uniform funding]\label{ex:track-record}
Two researchers, harmonic production, $A = 1$, budget $B = 2$. Researcher 1: $K_1 = 8$, $R_{10} = 8$ (productive and well resourced; expected output $\varphi_1(0) = 4$). Researcher 2: $K_2 = 4$, $R_{20} = 1$ (capable but with few resources; $\varphi_2(0) = 4/5$). Allocating in proportion to expected output gives $g = (5/3,\, 1/3)$ and total output $285/53 \approx 5.38$. Uniform funding, $g = (1, 1)$, gives $284/51 \approx 5.57$. The gap rule ($c = 3/4$) gives $g^* = (0, 2)$ and $40/7 \approx 5.71$.
\end{example}

\begin{example}[funding the under-resourced can produce less than uniform funding]\label{ex:scarcity}
Two researchers, harmonic production, $A = 1$, budget $B = 2$. Researcher 1: $K_1 = 10$, $R_{10} = 5$. Researcher 2: $K_2 = 1/2$, $R_{20} = 0$ (the scarcest resources in the population). Directing the budget to the researcher with the fewest resources gives total output $56/15 \approx 3.73$. Uniform funding gives $49/12 \approx 4.08$. The gap rule ($c = 2/3$) gives $g^* = (5/3,\, 1/3)$ and $21/5 = 4.2$.
\end{example}

% Figure (Appendix A): the two intuitive schemes on the population of Figure 1.
% Same 75 researchers as Figure 1 (fig1h-all.dat). Budget B_1 = 8.52, the total of the
% grants at the smaller budget in Figure 1 (c_1 = 2/3), divided equally among nine
% researchers (x = 0.947 each). Left: the nine with the highest expected output
% K R/(K + R) (fig1h-track.dat). Right: the nine with the fewest baseline resources
% (fig1h-under.dat). Expected output over no funding on this population (A = 1):
% gap rule +4.01, track record +1.98, under-resourced +2.09, uniform +1.99.
\begin{figure}[tp]
\centering
\begin{tikzpicture}
\begin{axis}[
  name=left,
  width=0.46\textwidth, height=5.2cm,
  axis lines*=left,
  xmin=0, xmax=8.6, ymin=0, ymax=7.6,
  xtick=\empty, ytick=\empty,
  x axis line style={-{Stealth[length=4pt]}, line width=0.4pt},
  y axis line style={-{Stealth[length=4pt]}, line width=0.4pt},
  xlabel={resources $R$},
  ylabel={capability $K$},
  ylabel style={font=\small}, xlabel style={font=\small},
  title={\small funding the track record},
  title style={yshift=-4pt},
  axis line style={line width=0.4pt},
  clip=false,
]
\addplot[only marks, mark=o, mark size=1.3pt, gray, line width=0.4pt] table[x=R, y=K] {fig1h-all.dat};
\addplot[quiver={u=\thisrow{u}, v=\thisrow{v}}, -{Stealth[length=3pt]}, gray, line width=0.5pt] table {fig1h-track.dat};
\addplot[only marks, mark=*, mark size=1.5pt, black] table[x=R, y=K] {fig1h-track.dat};
\end{axis}
\begin{axis}[
  at={(left.outer east)}, anchor=outer west, xshift=4pt,
  width=0.46\textwidth, height=5.2cm,
  axis lines*=left,
  xmin=0, xmax=8.6, ymin=0, ymax=7.6,
  xtick=\empty, ytick=\empty,
  x axis line style={-{Stealth[length=4pt]}, line width=0.4pt},
  y axis line style={-{Stealth[length=4pt]}, line width=0.4pt},
  xlabel={resources $R$},
  xlabel style={font=\small},
  title={\small funding the under-resourced},
  title style={yshift=-4pt},
  axis line style={line width=0.4pt},
  clip=false,
]
\addplot[only marks, mark=o, mark size=1.3pt, gray, line width=0.4pt] table[x=R, y=K] {fig1h-all.dat};
\addplot[quiver={u=\thisrow{u}, v=\thisrow{v}}, -{Stealth[length=3pt]}, gray, line width=0.5pt] table {fig1h-under.dat};
\addplot[only marks, mark=*, mark size=1.5pt, black] table[x=R, y=K] {fig1h-under.dat};
\end{axis}
\end{tikzpicture}
\caption{Two heuristics on the population of Fig.~\ref{fig:p1}. The budget of that figure's smaller frontier is divided equally among nine researchers. Left: funding the track record gives the grants (arrows) to the nine researchers with the highest expected research output, most of them well resourced and of middling capability. Right: funding the under-resourced gives the grants to the nine researchers with the fewest resources, most of them of low capability. On this population the gap rule adds 4.0 units of expected research output over no funding; funding the track record adds 2.0, funding the under-resourced 2.1, and dividing the budget uniformly across all 75 researchers 2.0.}
\label{fig:heuristics}
\end{figure}

